\subsection{Full Classification Examples}

\begin{examplebox}[A parameter in the coefficient matrix]
Classify
\[
\begin{cases}
x+y=3,\\
2x+ay=6.
\end{cases}
\]
The coefficient determinant is
\[
\det\begin{pmatrix}1&1\\2&a\end{pmatrix}=a-2.
\]
If $a\neq2$, there is exactly one solution. If $a=2$, the second equation is twice the first, so there are infinitely many solutions.
\end{examplebox}

\begin{examplebox}[A parameter in the right-hand side]
Classify
\[
\begin{cases}
x+y=3,\\
2x+2y=a.
\end{cases}
\]
The coefficient determinant is always zero. Comparing with twice the first equation gives:
\begin{itemize}
\item if $a=6$, the equations coincide and there are infinitely many solutions;
\item if $a\neq6$, the equations are inconsistent and there is no solution.
\end{itemize}
\end{examplebox}

\begin{interpretationbox}[Determinant first, compatibility second]
A parameter in the coefficient matrix may create isolated values where uniqueness fails. A parameter in the right-hand side may decide whether dependent equations agree or contradict one another.
\end{interpretationbox}
